\chapter{Sobre el uso de herramientas IA como chatGPT}
En la era digital actual, las herramientas basadas en Inteligencia Artificial (IA), como ChatGPT de OpenAI, se han vuelto cada vez más populares y accesibles. Estos sistemas ofrecen asistencia rápida y diversa en una amplia gama de tareas, desde la generación de texto hasta la resolución de consultas complejas. Sin embargo, es crucial entender la manera adecuada de utilizar estas herramientas para garantizar que la información generada sea confiable y útil.

\section{El rol de las herramientas IA: asistente, no autoridad}
Herramientas IA, como ChatGPT, deben considerarse un asistente en el proceso de generación de contenido, no una fuente autoritativa o final.  

A pesar de su capacidad avanzada para generar respuestas coherentes y, a menudo, bien informadas, estas herramientas se basan en la información disponible hasta una fecha de corte, lo cual puede resultar en la omisión de desarrollos recientes o cambios significativos en el conocimiento establecido.  

Además, debido a su naturaleza predictiva, estas herramientas pueden generar información no sustentada en hechos reales, introduciendo elementos inventados o especulativos en sus respuestas. 

\section{Responsabilidad personal en la creación de contenidos}
Cuando se utiliza una herramienta IA, es importante: 

Verificar la información proporcionada con fuentes confiables y actualizadas. No asumir que la respuesta es definitiva, especialmente en temas que cambian rápidamente o que requieren precisión técnica. 

Evaluar críticamente las respuestas. Analizar si son lógicas, coherentes y si se alinean con lo que ya se sabe sobre el tema. 

Utilizar las herramientas IA como un complemento a la investigación o proceso de escritura, no como la única fuente de información.

Aunque las herramientas IA pueden generar contenido basado en indicaciones, es esencial que cualquier trabajo presentado como propio refleje las ideas y comprensión originales del estudiante. 

\section{Mención de herramientas IA en documentos académicos}
En documentos académicos, es esencial aclarar el uso de herramientas IA como asistentes en el proceso de investigación o desarrollo de contenido, y evitar referenciarlas directamente como fuentes de información primaria.

En la descripción de la metodología, se debe indicar de manera precisa cómo y con qué propósito se emplearon estas herramientas IA resaltando, por ejemplo, su rol en la generación de ideas, estructuración preliminar de argumentos, soporte para el diseño o implementación de algoritmos, análisis inicial de datos, entre otros. Esta mención debe ir acompañada de una crítica reflexiva sobre la información obtenida, subrayando la responsabilidad del estudiante en la verificación de la veracidad y relevancia de los datos a través de fuentes confiables. Al hacerlo de esta forma, se reconoce el valor de la IA como complemento al esfuerzo humano, asegurando al mismo tiempo la rigurosidad y credibilidad del trabajo. 

Al mencionar herramientas IA en el documento, debemos citarlas y referenciarlas apropiadamaente. En la página de APA se sugiere hacerlo de esta forma: 

Cita en el texto:  

(OpenAI, 2024) 

Referencia:  

OpenAI. (2024). ChatGPT v3.5 [Large language model].   
https://chat.openai.com/chat. 